
\documentclass[a4paper,12pt]{article}
\usepackage[utf8]{inputenc}
\usepackage[english,slovene]{babel}
%\usepackage{blindtext}
%\usepackage{pdflscape}
\usepackage[tikz]{bclogo}
\usepackage{qrcode,marvosym}
\usepackage{pgf,tikz,pgfplots}
\pgfplotsset{compat=1.14}
\usepackage{mathrsfs}
\usetikzlibrary{arrows}
\usepackage{graphicx}
\newcommand{\ds}{\displaystyle}
\usepackage{fancybox}
\usepackage{amsfonts}
\usepackage{amsmath}
%\usepackage{color}
\usepackage{wallpaper}
\usepackage{hyperref}
\usepackage[cp1250]{inputenc}
\usepackage{array}%%%%%%%%%%%%%%%%%%%%%%%%%%%%%%%%%%%%%%%%%%%%
\newcolumntype{M}[1]{>{\centering\arraybackslash}m{#1}}%%%%
\newcolumntype{N}{@{}m{0pt}@{}}%%%%%%%%%%%%%%%%%%%%%%%%%%%%
   \usepackage{fancyhdr,enumerate}

 \newcommand{\ora}{\overrightarrow}
 \renewcommand{\headrulewidth}{3pt}
 \renewcommand{\headrule}{\hbox to\headwidth{%
   \color{gray}\leaders\hrule height \headrulewidth\hfill}}
    \renewcommand{\footrulewidth}{2pt}
    \renewcommand{\footrule}{\hbox to\headwidth{%
  \color{brown}\leaders\hrule height \footrulewidth\hfill}}
\usepackage{gensymb}
\usepackage{amssymb} 

  \usepackage{fancyhdr}
  \textwidth 20 true cm
  \oddsidemargin -2 true cm
  \evensidemargin -2 true cm
  \topmargin -3.5 true cm
  \textheight 27.5 true cm
  \linespread{1.2}
    \renewcommand{\headrulewidth}{0.3pt}
   \renewcommand{\footrulewidth}{1.9pt}
  \definecolor{qqwuqq}{rgb}{1,0,0 }
  \definecolor{qqqqff}{rgb}{0,0,1}
  \definecolor{qqffqq}{rgb}{0,1,0}
  \definecolor{ffqqqq}{rgb}{1,0,0}
\definecolor{zzuxux}{rgb}{0,0,0}%{0.92,0.03,0.03}
\usepackage{mdframed}
\usepackage{multicol}
 \definecolor{bgblue}{RGB}{0,0,253}
\usepackage{xcolor}
\usepackage{eurosym}
%%%%%%%%%%%%%%%%%%%%%%%%%%%%%%%%%%%%%%%%%%%%%%%%%%%%%%%%
\newcounter{tocke}
\setcounter{tocke}{0}
\usepackage{tikz-3dplot}
%%%%%%%%%%%%%%%%%%%%%%%%%%%%%%%%%%%%%%%%%%%%%%%%%%%%%%%%%%
\mdfdefinestyle{theoremstyle}{%
linecolor=brown!40,linewidth=1pt,%
frametitlerule=true,%
frametitlebackgroundcolor=brown!50,
innertopmargin=\topskip,
}

\mdfdefinestyle{theoremstyle1}{%
linecolor=brown,middlelinewidth=1pt,%
frametitlerule=true,%
apptotikzsetting={\tikzset{mdfframetitlebackground/.append style={%
shade,left color=black!40, right color=gray!10},
mdfbackground/.append style={
top color=white!100!white,
bottom color=black!5
},
}}, 
frametitlerulecolor=gray,
frametitlerulewidth=2pt,
innertopmargin=\topskip,
innerbottommargin=\topskip,
}
\mdtheorem[style=theoremstyle1]{naloga}{Naloga}
\mdtheorem[style=theoremstyle]{resitev}{Analiza Naloge}
\mdtheorem[style=theoremstyle]{kriterij}{\v{S}tevilo dose\v{z}enih to\v{c}k na testu}
%%%%%%%%%%%%%%%%%%%%%%%%%%%%%%%%%%%%%%%%%%%%%%%%%%%%%%%%%%
\usepackage[table]{xcolor}
\usepackage{titlesec}
\usepackage{venndiagram}
\usepackage[printwatermark]{xwatermark}
\usepackage{xcolor}
\newwatermark[allpages,color=brown!2,angle=45,scale=5,
xpos=-5,ypos=0]{matej.info}


%%%%%%%%%%%%%%%%%%%%%%%%%%%%%%%%%%%%%%%%%%%%%%%%%%%%%%%%%%
\titleformat{\subsection}
  {\normalfont\Large\bfseries\color{brown}}
  {\thesubsection}
  {1em}
  {}

\titleformat{\subsubsection}
  {\normalfont\large\bfseries\color{brown}}
  { \thesubsubsection}
  {1em}
  {}
%%%%%%%%%%%%%%%%%%%%
\pagestyle{fancy}
\renewcommand\headrulewidth{0pt}
\lhead{}\chead{}\rhead{}
\rfoot{\vspace*{-2\baselineskip}}


\renewcommand{\vec}{\overrightarrow}

%%%%%%%%%%%%%%%%%%%%%%55\Large


\usepackage{etoolbox}

\newcounter{secitemcount}      % števec itemov v sekciji
\newcounter{totalitemcount}    % skupni števec vseh nalog

\newcommand{\allresults}{}     % shranjevanje rezultatov

% ob začetku vsake sekcije resetiramo lokalni števec
\pretocmd{\section}{%
  \setcounter{secitemcount}{0}%
  \gdef\currentsectiontitle{}%
}{}{}

% prestrežemo naslov sekcije
\pretocmd{\section}{\gdef\currentsectiontitle{#1}}{}{}

% vsak \item poveča oba števca
\pretocmd{\item}{%
  \stepcounter{secitemcount}%
  \stepcounter{totalitemcount}%
}{}{}

% ob koncu enumerate avtomatsko shranimo rezultat
\AtEndEnvironment{enumerate}{%
  \gappto{\allresults}{%
    \resultentry{\currentsectiontitle}{\arabic{secitemcount}}%
  }%
}

% kako se izpiše posamezen rezultat v Rešitvah
\newcommand{\resultentry}[2]{%
  \subsection*{#1 (število nalog: #2)}%
}

%%%%%%%%%%%%%%%%%%%%%%%%%%%%%%%%%%%%%%%%%%%%%%%%%%%
% Definicija rjave barve
\definecolor{mybrown}{RGB}{140, 90, 40}

% SECTION – rjav okvir, bela notranjost
\titleformat{\section}
  {\large\bfseries}
  {%
    \fcolorbox{mybrown}{white}{\rule{1.2ex}{1.2ex}}%
    \quad\thesection
  }
  {0.8em}
  {}

% SUBSECTION – manjši rjav okvir, bela notranjost
\titleformat{\subsection}
  {\bfseries}
  {%
    \fcolorbox{mybrown}{white}{\rule{0.9ex}{0.9ex}}%
    \quad\thesubsection
  }
  {0.8em}
  {}

% --- Naslov dokumenta z rjavim kvadratkom ---
\newcommand{\DocTitle}[1]{%
  \begin{center}
    {\rule{\textwidth}{0.1mm}\\[0.4cm]
    \LARGE
    \fcolorbox{mybrown}{gray!40}{\rule{1.4ex}{1.4ex}}
    \quad #1
    }
  \end{center}
  \vspace{1em}
}
\newcounter{tocke}
\setcounter{tocke}{0}

\newmdenv[
    leftmargin=0pt,
    rightmargin=0pt,
    innerleftmargin=12pt,
    innerrightmargin=12pt,
    innertopmargin=2pt,
    innerbottommargin=8pt,
    skipabove=2pt,
    skipbelow=12pt,
    linewidth=20pt,
    linecolor=brown!60!black,
    topline=false,
    bottomline=false,
    rightline=false,
    backgroundcolor=gray!20,
    nobreak=false, % Omogoči prelom
    splittopskip=0pt,
    splitbottomskip=0pt
]{crta} 

\begin{document}

\pagestyle{fancy}
\lhead{}\chead{}\rhead{}

\begin{center}
\Huge {\textbf{Test – Odvod funkcije}}
\end{center}


% ===================== 1 =====================

\def\tockeA{3+4}
\begin{crta}
\subsection*{ Izračunaj odvod funkcije: \hfill{\tockeA}}
\addtocounter{tocke}{\tockeA}

\begin{enumerate}[a.)]
\item  $f(x)=3x^4-5x^2+2x-7$

\item $g(x)=\dfrac{3}{x^2}$
\end{enumerate}
\end{crta}
% ===================== 2 =====================

\def\tockeB{2+2}
\addtocounter{tocke}{\tockeB}
\subsection{ Izračunaj odvod v dani točki:}
\begin{enumerate}[a.)]
\item  $f(x)=x^3-2x^2+x$, \quad $x_0=2$\item  $g(x)=\dfrac{2x+1}{x-1}$, \quad $x_0=3$
\end{enumerate}


\vspace{1cm}

% ===================== 3 =====================
\def\tockeC{2+2+2}
\addtocounter{tocke}{\tockeC}
\subsection{ Določi enačbo tangente na graf funkcije v dani točki:} \hfill {tockeC}

a) $f(x)=x^2$, \quad $x_0=1$

b) $g(x)=\sqrt{x}$, \quad $x_0=4$

c) $h(x)=\frac{1}{x}$, \quad $x_0=1$

\vspace{1cm}

% ===================== 4 =====================
\def\tockeD{2+2+2}
\addtocounter{tocke}{\tockeD}
\subsection{Določi kot, ki ga tangenta na graf funkcije v dani točki oklepa z osjo $x$:}

a) $f(x)=\frac{3}{x}$, \quad $x_0=1$

b) $g(x)=\dfrac{1}{3}x^3$, \quad $x_0=2$


\vspace{1cm}

% ===================== 5 =====================
\def\tockeE{3+3}
\addtocounter{tocke}{\tockeE}
\subsection{Izračunaj kot med krivuljama v presečišču:}

a) $f(x)=x^2$, \quad $g(x)=x^2-3x$

b) $f(x)=x^2+1$, \quad $g(x)=-x^2+3$

\vfill


\newpage

% ===================== REŠITVE =====================
\section*{Rešitve}

\subsection*{1}
a)
\[
f'(x)=12x^3-10x+2
\]

b)
\[
g(x)=3x^{-2} \Rightarrow g'(x)=-6x^{-3}=-\frac{6}{x^3}
\]

\subsection*{2}
a)
\[
f'(x)=3x^2-4x+1
\]
\[
f'(2)=12-8+1=5
\]

b)
\[
g'(x)=\frac{2(x-1)-(2x+1)}{(x-1)^2}
=\frac{2x-2-2x-1}{(x-1)^2}
=\frac{-3}{(x-1)^2}
\]
\[
g'(3)=\frac{-3}{4}
\]

\subsection*{3}
Uporabimo: $y=f'(x_0)(x-x_0)+f(x_0)$

a)
\[
f'(x)=2x,\quad f'(1)=2,\quad f(1)=1
\]
\[
y=2(x-1)+1=2x-1
\]

b)
\[
g'(x)=\frac{1}{2\sqrt{x}},\quad g'(4)=\frac{1}{4},\quad g(4)=2
\]
\[
y=\frac{1}{4}(x-4)+2
\]

c)
\[
h'(x)=-\frac{1}{x^2},\quad h'(1)=-1,\quad h(1)=1
\]
\[
y=-(x-1)+1=-x+2
\]

\subsection*{4}
Velja: $\tan \alpha = f'(x_0)$

a)
\[
f'(x)=-\frac{3}{x^2},\quad f'(1)=-3
\]
\[
\alpha=\arctan(-3)
\]

b)
\[
g'(x)=x^2,\quad g'(2)=4
\]
\[
\alpha=\arctan(4)
\]

\subsection*{5}
Formula:
\[
\tan \varphi = \left|\frac{k_1-k_2}{1+k_1k_2}\right|
\]

a)
Presečišče:
\[
x^2=x^2-3x \Rightarrow x=0
\]
\[
f'(x)=2x,\quad g'(x)=2x-3
\]
\[
k_1=0,\quad k_2=-3
\]
\[
\tan \varphi = 3
\]

b)
Presečišče:
\[
x^2+1=-x^2+3 \Rightarrow 2x^2=2 \Rightarrow x=\pm1
\]

Za $x=1$:
\[
f'(x)=2x=2,\quad g'(x)=-2x=-2
\]
\[
\tan \varphi = \left|\frac{2-(-2)}{1-4}\right|=\frac{4}{3}
\]

(enak rezultat dobimo tudi za $x=-1$)

\end{document}